\section{Introduction}
\label{sec:intro}
Cloud computing has become core infrastructure for public administrations and critical sections, but US is the main supplier of the most advanced cloud technology. The US CLOUD Act allows US authorities to demand data from American cloud providers, even when the data is stored outside the United States \parencite{doj2019cloudact}. For European governments handling sensitive data, this creates a conflict between technological capability and control. S3NS was created as a French answer to this problem. It is a company under French law, fully owned by Thales, that runs a trusted cloud based on Google Cloud technology \parencite{s3ns2025premi3ns}.

S3NS describes itself as an alliance between Thales, the French leader in
cybersecurity in Europe, and Google Cloud.
The two partners play different roles. S3NS is a French-law company fully
controlled by Thales, which runs the service in France, while Google Cloud
supplies the underlying technology, including Compute Engine, Cloud Storage,
Cloud SQL, Google Kubernetes Engine and BigQuery.
Every Google technology update is first received in a quarantine zone, where
S3NS analyses and validates it before it is deployed in an environment
administered only by S3NS \parencite{s3ns2025premi3ns}. Google therefore holds
no control over the operation, yet the service cannot function without its
technology. The resulting offering, PREMI3NS, obtained the SecNumCloud 3.2
qualification from the French cybersecurity agency ANSSI in December 2025, and
is used by around thirty early customers, including MGEN, Matmut, Club Med and
Birdz, a Veolia subsidiary. EDF has also announced that it has selected S3NS
\parencite{s3ns2025premi3ns}.

Demand for such services is shaped by regulation. French rules require
sensitive state data hosted with external providers to use SecNumCloud-qualified
services, a principle that began as government cloud policy and was written into
law through the SREN law, with implementing rules adopted in 2026
\parencite{decretsren2026}. S3NS competes with Bleu, a trusted cloud
built by Orange and Capgemini on Microsoft technology \parencite{lemagit2023usf},
and with independent French providers such as OVHcloud. Its ecosystem of
software vendors is still at an early stage \parencite{futurum2026}. The case
therefore raises a central tension: S3NS promises sovereignty, yet depends on a
partner whose technology it cannot replace.

This report asks how the division of roles between Thales and Google shapes
ecosystem dependencies, lock-in, platform pricing, security incentives and value
capture in S3NS. Section~\ref{sec:task1} to Section~\ref{sec:task5} analyse
these questions using the frameworks of \textcite{overby2021} and the course,
and Section~\ref{sec:conclusion} concludes. AI use follows the course
requirements and is documented in Appendices~A and~B.

